\section{แนวคิดพื้นฐานด้านการตรวจสอบสภาพผิวถนน}

\subsection{ความสำคัญของคุณภาพผิวถนน}

คุณภาพผิวถนนเป็นปัจจัยสำคัญที่ส่งผลโดยตรงต่อความปลอดภัยในการเดินทาง โดยเฉพาะแรงเสียดทานของผิวทางและความต้านทานการลื่นไถล ซึ่งมีผลต่อความสามารถของยานพาหนะในการควบคุมทิศทางและหยุดรถในสถานการณ์ฉุกเฉิน ผิวถนนที่เสื่อมสภาพหรือมีแรงเสียดทานต่ำอาจเพิ่มความเสี่ยงต่อการสูญเสียการควบคุมยานพาหนะ โดยเฉพาะในสภาวะถนนเปียก องค์การอนามัยโลก (World Health Organization: WHO) ระบุว่า คุณภาพของโครงสร้างพื้นฐานทางถนน รวมถึงการออกแบบและการบำรุงรักษา เป็นปัจจัยสำคัญที่ส่งผลต่อความเสี่ยงของการบาดเจ็บและการเสียชีวิตจากอุบัติเหตุทางถนน \parencite{who2023} ในทำนองเดียวกัน สำนักงานทางหลวงแห่งสหรัฐอเมริกา (Federal Highway Administration: FHWA) จัดให้การบริหารจัดการแรงเสียดทานของผิวทางเป็นหนึ่งในมาตรการด้านความปลอดภัยที่ได้รับการพิสูจน์แล้วว่าสามารถช่วยลดอุบัติเหตุได้อย่างมีประสิทธิภาพ \parencite{fhwa2023friction}

นอกจากผลกระทบด้านความปลอดภัยแล้ว คุณภาพผิวถนนยังส่งผลโดยตรงต่อต้นทุนการดำเนินงานของยานพาหนะ ถนนที่มีความขรุขระสูงทำให้การใช้เชื้อเพลิงเพิ่มขึ้น ส่งผลให้ยาง ระบบกันสะเทือน และส่วนประกอบของยานพาหนะเกิดการสึกหรอเร็วขึ้น และนำไปสู่ค่าบำรุงรักษาที่สูงขึ้น FHWA ระบุว่า ความเรียบของผิวทางมีผลต่อการใช้เชื้อเพลิงและประสิทธิภาพของการขนส่งโดยรวม ดังนั้น การบำรุงรักษาผิวถนนให้อยู่ในสภาพดีจึงมีส่วนช่วยลดต้นทุนการขนส่งในระยะยาวได้ \parencite{fhwa2016}

ในมุมมองของการบริหารจัดการสินทรัพย์โครงสร้างพื้นฐาน คุณภาพผิวถนนถือเป็นตัวชี้วัดสำคัญสำหรับการตัดสินใจด้านงบประมาณและการวางแผนบำรุงรักษา ตัวชี้วัด เช่น ความขรุขระของผิวถนน ร่องล้อ และรอยแตกร้าวของผิวทาง ถูกนำมาใช้เพื่อประเมินว่าถนนควรได้รับการบำรุงรักษาเชิงป้องกัน หรือควรเข้าสู่กระบวนการฟื้นฟูสภาพถนนในระดับที่มากขึ้น สมาคมถนนโลก (World Road Association: PIARC) ชี้ว่า การรักษาคุณภาพผิวทางให้อยู่ในระดับที่เหมาะสมสามารถช่วยลดต้นทุนตลอดอายุการใช้งานของโครงสร้างพื้นฐานได้ เนื่องจากการแก้ไขความเสียหายในระยะเริ่มต้นมีต้นทุนต่ำกว่าการปล่อยให้ความเสียหายสะสมจนต้องบูรณะขนาดใหญ่ \parencite{piarcLifecycle} ด้วยเหตุนี้ การตรวจสอบสภาพผิวถนนอย่างต่อเนื่องและทันเวลาจึงเป็นองค์ประกอบสำคัญของการบริหารจัดการโครงสร้างพื้นฐานทางถนนอย่างมีประสิทธิภาพ

\subsection{ดัชนี IRI และมาตรฐานการวัดสภาพผิวถนน}

ดัชนีความขรุขระสากล (International Roughness Index: IRI) เป็นตัวชี้วัดมาตรฐานที่ใช้ประเมินความขรุขระของผิวถนนตามแนวการเคลื่อนที่ของยานพาหนะ โดยมีหน่วยเป็นเมตรต่อกิโลเมตร ค่า IRI สะท้อนระดับความไม่เรียบสะสมของผิวถนนต่อระยะทาง 1 กิโลเมตร ดัชนีดังกล่าวพัฒนาขึ้นจากแบบจำลองรถหนึ่งในสี่ส่วน เพื่อจำลองการตอบสนองของระบบกันสะเทือนของยานพาหนะต่อโปรไฟล์ผิวถนน ทำให้ IRI เป็นตัวชี้วัดที่มีความสัมพันธ์กับประสบการณ์การขับขี่จริงและผลกระทบที่เกิดขึ้นกับยานพาหนะ \parencite{sayers1986roughness} ด้วยเหตุนี้ ค่า IRI จึงถูกนำมาใช้อย่างแพร่หลายในระบบบริหารจัดการผิวทาง เพื่อประเมินระดับคุณภาพของผิวถนนและสนับสนุนการตัดสินใจด้านการบำรุงรักษา

ค่า IRI มีความสัมพันธ์โดยตรงกับต้นทุนการดำเนินงานของยานพาหนะ เนื่องจากถนนที่มีค่า IRI สูงทำให้เกิดแรงสั่นสะเทือนและแรงต้านการเคลื่อนที่เพิ่มขึ้น ส่งผลให้การใช้เชื้อเพลิงเพิ่มขึ้น ชิ้นส่วนของยานพาหนะสึกหรอเร็วขึ้น และค่าบำรุงรักษาโดยรวมสูงขึ้น แบบจำลองการพัฒนาและบริหารจัดการทางหลวง (Highway Development and Management: HDM) ของธนาคารโลกแสดงให้เห็นว่า ความขรุขระของถนนเป็นปัจจัยสำคัญที่ส่งผลต่อต้นทุนการดำเนินงานของยานพาหนะ และมีผลกระทบเชิงเศรษฐกิจต่อระบบขนส่งโดยรวม \parencite{sayers1986roughness} แม้ว่า IRI ไม่ได้วัดแรงเสียดทานของผิวทางโดยตรง แต่ถนนที่มีความขรุขระสูงมักมีความสัมพันธ์กับการเสื่อมสภาพของผิวทาง การเกิดร่องล้อ และความไม่เรียบของพื้นผิว ซึ่งส่งผลต่อเสถียรภาพของยานพาหนะและอาจเพิ่มความเสี่ยงต่อการเกิดอุบัติเหตุ \parencite{fhwa2016}

ในงานวิจัยด้านการตรวจสอบสภาพถนน ค่า IRI มักถูกใช้เป็นตัวชี้วัดอ้างอิงสำหรับประเมินระดับความขรุขระของผิวทาง เนื่องจากมีนิยามทางวิศวกรรมที่ชัดเจนและสามารถใช้เปรียบเทียบผลการวัดระหว่างวิธีการต่าง ๆ ได้ อย่างไรก็ตาม การได้มาซึ่งค่า IRI ที่มีความแม่นยำต้องอาศัยเครื่องวัดโปรไฟล์ผิวถนนหรืออุปกรณ์เฉพาะทางที่มีต้นทุนสูง และต้องดำเนินการโดยผู้ปฏิบัติงานที่มีความชำนาญ จึงไม่เหมาะต่อการขยายผลสู่การเก็บข้อมูลแบบระดมข้อมูลจากมวลชน ซึ่งอาศัยผู้ใช้ทั่วไปจำนวนมากในสภาพแวดล้อมจริง

ด้วยข้อจำกัดดังกล่าว งานวิจัยนี้จึงออกแบบการเก็บข้อมูลโดยใช้หน่วยวัดความเฉื่อย (Inertial Measurement Unit: IMU) และระบบระบุตำแหน่งบนพื้นโลก (Global Positioning System: GPS) บนสมาร์ตโฟน ร่วมกับป้ายกำกับเชิงกลุ่มที่กำหนดจากหลักฐานที่สังเกตได้โดยตรง ได้แก่ ภาพวิดีโอจากกล้องบันทึกหน้ารถและสัญญาณแรงสั่นสะเทือนที่เกิดขึ้นอย่างชัดเจน แนวทางนี้เหมาะสมกับบริบทของการระดมข้อมูลจากมวลชน เนื่องจากป้ายกำกับเชิงกลุ่มสามารถรวบรวมและตรวจสอบข้ามเที่ยวการเดินทางจากผู้ใช้หลายรายในเส้นทางเดียวกันได้ โดยไม่จำเป็นต้องพึ่งพาค่าการวัดแบบสัมบูรณ์ในระดับเดียวกับค่า IRI

\subsection{วิธีการดั้งเดิมในการตรวจสอบสภาพผิวถนน}

วิธีการดั้งเดิมในการวัดโปรไฟล์ผิวถนนสามารถจำแนกได้เป็น 3 ประเภทหลักตามหลักการวัด ได้แก่ วิธี Rod and Level วิธี Profilograph และวิธี Inertial Profiler โดยแต่ละวิธีมีระดับความแม่นยำ ต้นทุน และความเหมาะสมต่อการใช้งานที่แตกต่างกัน

วิธีแรกคือ Rod and Level ซึ่งเป็นวิธีการวัดแบบสถิต โดยใช้กล้องระดับความแม่นยำสูงร่วมกับไม้บรรทัดวัดระดับ เพื่อสำรวจค่าความสูงสัมพัทธ์ของผิวถนนทีละจุดตลอดแนวล้อของยานพาหนะ โดยทั่วไปมีการเก็บค่าที่ระยะห่างประมาณ 0.3 เมตร วิธีนี้จัดอยู่ในระดับ Class I ตามมาตรฐานของ American Society for Testing and Materials หรือ ASTM และมักใช้เป็นวิธีอ้างอิงสำหรับสอบเทียบอุปกรณ์วัดชนิดอื่น เนื่องจากสามารถให้ความละเอียดในแนวดิ่งได้สูง ประมาณ 0.125–2.0 มิลลิเมตร ขึ้นอยู่กับระดับความขรุขระของผิวถนน \parencite{fhwa2016} อย่างไรก็ตาม วิธีนี้ต้องใช้ผู้ปฏิบัติงานหลายคน ใช้เวลานาน และมีอัตราการสำรวจต่ำ จึงไม่เหมาะต่อการสำรวจโครงข่ายถนนขนาดใหญ่

วิธีที่สองคือ Profilograph ซึ่งเป็นอุปกรณ์วัดแบบใช้โครงอ้างอิงขณะเคลื่อนที่ โดยใช้โครงความยาวประมาณ 24.75–25 ฟุตรองรับด้วยล้อหลายล้อ เพื่อวัดการเคลื่อนที่ในแนวดิ่งของล้อกลางเทียบกับโครงอ้างอิง ผลลัพธ์ที่ได้มักแสดงในรูปของดัชนีโปรไฟล์ ไม่ใช่โปรไฟล์ผิวถนนจริงโดยตรง เนื่องจากการตอบสนองต่อความถี่ของระบบขึ้นอยู่กับความยาวของโครงอุปกรณ์ ทำให้อุปกรณ์อาจขยายหรือลดทอนสัญญาณในบางช่วงความยาวคลื่นได้ วิธี Profilograph จึงเหมาะสำหรับการตรวจรับงานก่อสร้างและการระบุความผิดปกติเฉพาะจุด เช่น จุดนูนหรือจุดสะดุดของผิวทาง แต่ไม่เหมาะสำหรับการคำนวณค่า IRI หรือการสำรวจสภาพถนนในระดับโครงข่ายขนาดใหญ่ \parencite{fhwa2016}

วิธีที่สามคือ Inertial Profiler ซึ่งเป็นอุปกรณ์มาตรฐานที่ใช้กันอย่างแพร่หลายในการสำรวจโครงข่ายถนนในปัจจุบัน ระบบนี้ประกอบด้วยองค์ประกอบหลัก 3 ส่วน ได้แก่ มาตรวัดความเร่งสำหรับวัดการเคลื่อนที่ของตัวรถ เซ็นเซอร์วัดระยะด้วยเลเซอร์แบบไม่สัมผัสสำหรับวัดระยะจากตัวรถถึงผิวถนน และเครื่องมือวัดระยะทางสำหรับกำหนดตำแหน่งตามแนวยาวของถนน ระบบดังกล่าวสามารถทำงานได้ในขณะรถเคลื่อนที่ด้วยความเร็วสูง และเหมาะสำหรับการคำนวณค่า IRI ในการสำรวจถนนระดับโครงข่าย อย่างไรก็ตาม อุปกรณ์ประเภทนี้มีต้นทุนสูง ต้องได้รับการสอบเทียบอย่างสม่ำเสมอ และต้องอาศัยผู้ปฏิบัติงานที่มีทักษะเฉพาะทาง

โดยสรุป วิธี Rod and Level มีความแม่นยำสูงและเหมาะสำหรับใช้เป็นวิธีอ้างอิง แต่มีข้อจำกัดด้านเวลาและแรงงาน วิธี Profilograph เหมาะสำหรับการตรวจรับงานก่อสร้างและการตรวจหาความผิดปกติเฉพาะจุด ส่วนวิธี Inertial Profiler เหมาะสำหรับการสำรวจถนนในระดับโครงข่าย แต่ยังมีข้อจำกัดด้านต้นทุน การสอบเทียบ และความต้องการบุคลากรเฉพาะทาง ข้อจำกัดเหล่านี้เป็นเหตุผลสำคัญที่ทำให้การพัฒนาแนวทางการตรวจสอบสภาพถนนด้วยสมาร์ตโฟนและเซ็นเซอร์ต้นทุนต่ำมีความสำคัญต่อการขยายผลในสภาพแวดล้อมจริง ตารางที่ 2.1 สรุปการเปรียบเทียบวิธีการดั้งเดิมทั้ง 3 ประเภทในด้านหลักการวัด ความแม่นยำ ความเร็วสำรวจ และต้นทุน

\begingroup
\setlength{\tabcolsep}{3pt}
\renewcommand{\arraystretch}{1.20}
\setlength{\LTleft}{0pt}
\setlength{\LTright}{0pt}
\begin{longtable}{>{\raggedright\arraybackslash}p{0.15\textwidth}>{\raggedright\arraybackslash}p{0.34\textwidth}>{\raggedright\arraybackslash}p{0.18\textwidth}>{\raggedright\arraybackslash}p{0.25\textwidth}}
\caption{สรุปการเปรียบเทียบวิธีการดั้งเดิมในการตรวจสอบสภาพผิวถนน}\label{tab:traditional-road-methods}\\
\toprule
\textbf{วิธีการ} & \textbf{หลักการวัดและความแม่นยำ} & \textbf{ความเร็วและต้นทุน} & \textbf{การใช้งานหลักและข้อสังเกต} \\
\midrule
\endfirsthead
\SUTContinuedTableHeader{4}
\toprule
\textbf{วิธีการ} & \textbf{หลักการวัดและความแม่นยำ} & \textbf{ความเร็วและต้นทุน} & \textbf{การใช้งานหลักและข้อสังเกต} \\
\midrule
\endhead
\bottomrule
\endfoot
Rod and Level & Static surveying ทีละจุด ทุก 0.3 m; ASTM Class I; resolution 0.125–2.0 mm & เดิน (\textless{}10 วินาที/จุด); ต้นทุนอุปกรณ์ต่ำแต่ใช้แรงงานสูง & ใช้สอบเทียบอุปกรณ์อื่นและงานวิจัยระยะสั้น; ไม่เหมาะกับการสำรวจโครงข่ายขนาดใหญ่ \\
Profilograph & Rolling-reference beam ประมาณ 25 ฟุต; รายงาน Profile Index (PI) และไม่ใช่ true IRI & ประมาณ 2–3 mph; ต้นทุนต่ำถึงปานกลาง & ใช้ตรวจสอบคุณภาพงานก่อสร้างและ legacy PI specifications; ไม่เหมาะกับ network-level IRI survey \\
Inertial Profiler & Accelerometer + laser + DMI; IRI ratio 0.95–1.05 เทียบ reference และมี repeatability สูง & 80–88 km/h; ต้นทุนสูง & เหมาะกับ network-level IRI survey และการบริหารสินทรัพย์ แต่ต้องสอบเทียบและใช้บุคลากรเฉพาะทาง \\
\end{longtable}
\endgroup
\SUTTableNote{ดัดแปลงจาก \fullcite{fhwa2016} และ \fullcite{zang2018}}

โดยสรุป วิธีการวัดโปรไฟล์ผิวถนนแบบดั้งเดิมมีข้อจำกัดร่วมกันด้านต้นทุน ความซับซ้อนของอุปกรณ์ ความต้องการบุคลากรที่มีความชำนาญ และความสามารถในการขยายผลสู่การสำรวจถนนในวงกว้าง โดยเฉพาะในประเทศกำลังพัฒนาหรือพื้นที่ที่มีข้อจำกัดด้านงบประมาณ วิธีการที่ต้องอาศัยเครื่องมือเฉพาะทาง เช่น เครื่องวัดโปรไฟล์ผิวถนนแบบเฉื่อย จึงอาจไม่เหมาะต่อการติดตามสภาพถนนอย่างต่อเนื่องในระยะยาว งานทบทวนของ \textcite{alqaydi2024} ชี้ให้เห็นว่า แนวทางการใช้สมาร์ตโฟนร่วมกับเซ็นเซอร์ภายในอุปกรณ์เป็นทางเลือกที่มีต้นทุนต่ำและสามารถขยายผลได้มากกว่าวิธีสำรวจถนนแบบดั้งเดิม
